In this work, we introduced EliteDesign, an RNA design framework that uses the
gradient-free cross-entropy method to find high-probability RNA sequences for a given target structure.
Our method beats the state-of-the-art SamplingDesign on Eterna100 in both
speed and quality, and we show large gains over SAMFEO on ten 16S
rRNA targets while remaining much faster per puzzle. We further show that this method addresses a gap in current RNA design
methods by enabling the design of sequences longer than 1,000 nucleotides with relatively high
Boltzmann probability compared with other methods. We also show that our
method achieves higher Boltzmann probability on other RNA design datasets such as Rfam~\citep{taneda2011,kalvari2021}
and RNAsolo~\citep{adamczyk2022}.

\subsection{Why EliteDesign works}\label{sec:why}
EliteDesign keeps the backbone of SamplingDesign by searching for a high-probability sequence with a distribution over the design space. However, it changes how this distribution learns. SamplingDesign averages a gradient over all of its samples and takes a small step, whereas EliteDesign learns only from the best samples and moves halfway to them in one update. As a result, it needs far fewer scoring calls to LinearPartition. For example, on Eterna100 with $N=500$, EliteDesign samples about 61,000 sequences on average and takes about 120 rounds, whereas SamplingDesign samples about 4.35 million sequences and takes about 1,740 rounds.

EliteDesign also takes advantage of optimizing for two objectives, as shown in Eq.~\ref{eq:pair}, where either objective can be replaced without changing the algorithm. Since neither needs a gradient, even an objective that is not differentiable can be used.

\subsection{Limitations}\label{sec:limitations}
Despite these gains, EliteDesign has some limitations.
\begin{enumerate}
\item EliteDesign uses the same initialization strategy for every target, which may not work
  equally well for all structures. During optimization, it can concentrate too early on
  a narrow set of sequences and miss better solutions. This may explain the small losses to SamplingDesign on some short puzzles.
\item EliteDesign uses LinearPartition to estimate Boltzmann probabilities during search. Errors in these estimates
  can cause it to favor weaker candidates and miss better ones.
  Exact scoring at the end checks the final candidates but cannot recover sequences
  missed during search.
\item Tianshuo's Ceiling shows that room remains on many targets, for example on Snowflake 4, where EliteDesign reaches 47\% of the ceiling. However, a large gap can come either from the design or from a loose bound, so it does not prove that a better sequence exists.
\end{enumerate}

\subsection{Ablation study}\label{sec:ablation}
\textbf{TODO.}

\subsection{Future work}\label{sec:future}
We see several directions for extending this work.
\begin{enumerate}
\item Future work could add a local refinement stage similar to SAMFEO~\citep{zhou2023},
which might further improve NED and Boltzmann probability on some structures.
\item Because the objectives are interchangeable, EliteDesign can be applied to other objectives and to other design problems such as mRNA design~\citep{zhang2023}.
\item While EliteDesign does well on computational benchmarks, where designs are only scored with a thermodynamic model, the designs still need to be tested in the laboratory.
\end{enumerate}

\subsection{Code availability}\label{sec:code}
Our code is available at \url{https://github.com/JSONFN21/EliteDesign}.
